% ---------------------------------------------------------------------
\chapter{Meldungen verstehen}
\label{kap:meldungen}

\begin{lernziele}
\begin{itemize}[nosep]
  \item die häufigsten Meldungen von Git und \cmd{gh} einordnen,
  \item zwischen harmlosen Hinweisen und echten Fehlern unterscheiden.
\end{itemize}
\end{lernziele}

Git unterscheidet drei Arten von Meldungen. Sie beginnen meist mit einem Schlüsselwort:

\begin{center}
\renewcommand{\arraystretch}{1.25}
\begin{tabularx}{\linewidth}{@{}>{\ttfamily}l>{\raggedright\arraybackslash}X@{}}
  \toprule
  \normalfont\textbf{Beginnt mit} & \textbf{Bedeutung} \\
  \midrule
  hint:    & Tipp; der Befehl war erfolgreich oder Git schlägt den nächsten Schritt vor \\
  warning: & Hinweis; der Befehl wurde ausgeführt, aber etwas ist bemerkenswert \\
  error: / fatal: & Fehler; der Befehl wurde \textbf{nicht} ausgeführt \\
  \bottomrule
\end{tabularx}
\end{center}

\section{Zustandsmeldungen}

\meldung{nothing to commit, working tree clean}
  {Es gibt keine Änderungen. Arbeitsordner und letzter Commit sind identisch.}
  {Nichts. Alles ist gespeichert.}

\meldung{Your branch is ahead of \textquotesingle{}origin/main\textquotesingle{} by 1 commit.}
  {Ein lokaler Commit ist noch nicht auf GitHub.}
  {\cmd{git push}}

\meldung{Your branch is behind \textquotesingle{}origin/main\textquotesingle{} by 2 commits.}
  {Auf GitHub gibt es neue Commits, die lokal fehlen.}
  {\cmd{git pull}}

\section{Hinweise}

\meldung{warning: in the working copy of \textquotesingle{}datei.py\textquotesingle, LF will be replaced by CRLF}
  {Die Datei verwendet Zeilenenden im Linux-Format. Git wird sie unter Windows anpassen
   (siehe Abschnitt~\ref{sec:zeilenenden}).}
  {Nichts. Reine Information.}

\section{Fehler}

\meldung{error: src refspec ads does not match any}
  {\cmd{git push} soll den Branch \cmd{ads} hochladen, der lokal nicht existiert.
   Meist wurde \cmd{git switch -c ads} vergessen und auf \cmd{main} committet.}
  {Mit \cmd{git branch} prüfen. Den Branch nachträglich anlegen: \cmd{git switch -c ads},
   dann erneut pushen.}

\meldung{no pull requests found for branch "ads"}
  {\cmd{gh pr merge} findet keinen Pull Request, weil \cmd{gh pr create} übersprungen wurde.}
  {Zuerst \cmd{gh pr create --fill}, dann mergen.}

\meldung{fatal: not a git repository (or any of the parent directories): .git}
  {Das Terminal steht in einem Ordner, der zu keinem Repository gehört.}
  {Mit \cmd{cd} in den Repository-Ordner wechseln.}

\meldung{! [rejected] main -> main (fetch first)}
  {Auf GitHub gibt es Commits, die lokal fehlen. Git verweigert das Überschreiben, um
   fremde Arbeit zu schützen.}
  {Zuerst \cmd{git pull}, gegebenenfalls Konflikte lösen, dann \cmd{git push}.}

\meldung{CONFLICT (content): Merge conflict in README.md}
  {Dieselbe Stelle wurde auf zwei Wegen unterschiedlich geändert.}
  {Siehe Kapitel~\ref{kap:zusammenarbeit}.}

\meldung{could not prompt: Incorrect function. You appear to be running in MinTTY}
  {Git Bash kann die Auswahlmenüs von \cmd{gh} nicht darstellen.}
  {Alle Antworten als Optionen mitgeben oder PowerShell verwenden.}

\meldung{refusing to allow an OAuth App to create or update workflow \textasciigrave.github/workflows/ci.yml\textasciigrave{} without \textasciigrave workflow\textasciigrave{} scope}
  {Der Push enthält eine Datei unter \cmd{.github/workflows/}. Die Anmeldung von \cmd{gh}
   hat dafür keine Berechtigung, weil Workflows automatisch Code ausführen.}
  {\cmd{gh auth refresh -h github.com -s workflow} und im Browser bestätigen.}

\section{Zeilenenden}
\label{sec:zeilenenden}

Ein Zeilenumbruch ist in einer Textdatei ein unsichtbares Zeichen. Linux und macOS
verwenden eines, \emph{Line Feed} (LF). Windows verwendet zwei, \emph{Carriage Return} und
\emph{Line Feed} (CRLF). Die Namen stammen von der Schreibmaschine: Wagen zurückschieben,
Papier eine Zeile weiterdrehen.

Mit der Einstellung \cmd{core.autocrlf = true} wandelt Git beim Auschecken unter Windows in
CRLF um und speichert im Repository LF. So passt jede Datei auf jedem System, und
Unterschiede in den Zeilenenden erscheinen nicht als Änderungen.

\begin{zusammenfassung}
\begin{itemize}[nosep]
  \item \cmd{hint:} und \cmd{warning:} sind Informationen, \cmd{error:} und \cmd{fatal:} echte Fehler.
  \item Fast jede Fehlermeldung nennt Ursache und Lösung. In Ruhe lesen lohnt sich.
\end{itemize}
\end{zusammenfassung}
